
\documentclass{article}
\usepackage{colortbl}
\usepackage{makecell}
\usepackage{multirow}
\usepackage{supertabular}

\begin{document}

\newcounter{utterance}

\centering \large Interaction Transcript for game `eqbench', experiment `eq\_bench', episode 58 with reflexion{-}Qwen3.5{-}27B.
\vspace{24pt}

{ \footnotesize  \setcounter{utterance}{1}
\setlength{\tabcolsep}{0pt}
\begin{supertabular}{c@{$\;$}|p{.15\linewidth}@{}p{.15\linewidth}p{.15\linewidth}p{.15\linewidth}p{.15\linewidth}p{.15\linewidth}}
    \# & \multicolumn{2}{c}{Player} && \multicolumn{2}{c}{Game Master} \\
    \hline

    \theutterance \stepcounter{utterance}  
    & & & \multicolumn{4}{p{0.6\linewidth}}{
        \cellcolor[rgb]{0.9,0.9,0.9}{
            \makecell[{{p{\linewidth}}}]{
                \texttt{\tiny{[P1$\langle$GM]}}
                \texttt{You have attempted this same task before and did not succeed. These are the plans you wrote after those attempts.} \\
\texttt{Trial \#1: New plan: My previous attempt failed because I assigned a high score to "Guilt" (9) and a non{-}zero score to "Determination" (3), which likely contradicted the specific emotional nuance of the dialogue where Claudia's primary feeling is defensive justification rather than deep remorse, and her resolve is overshadowed by the conflict. Tyler's accusations ("afterthought," "burden") and Claudia's defensive replies ("That's not fair," "I had to") suggest she feels **Resentment** towards his unfair characterization of her actions more than she feels Guilt. Additionally, the prompt implies a single dominant emotion or a specific distribution that I missed. Next time, I will re{-}evaluate the dialogue to identify the *primary* emotional driver: Claudia is defending her past choices against perceived unfairness, which points to **Resentment** as the highest score (e.g., 8{-}9), while **Guilt** should be lower (e.g., 2{-}3) as she explicitly denies leaving him behind. **Relief** should remain 0 as the conversation is tense, and **Determination** should be low or 0 as she is reacting defensively rather than proactively. I will adjust the scores to reflect this hierarchy: Resentment > Guilt > Determination > Relief.} \\
\\ 
\\ 
\\ 
\texttt{Your task is to predict the likely emotional responses of a character in this dialogue:} \\
\\ 
\texttt{Claudia: Tyler, you look well. Teaching suits you.} \\
\texttt{Tyler: Yeah, well, someone's gotta do the dirty work.} \\
\texttt{Claudia: I didn't mean it like that. I...} \\
\texttt{Tyler: I know what you meant. You always do, don't you?} \\
\texttt{Claudia: Not always, no. But I try to understand, at least.} \\
\texttt{Tyler: Understand? Like you understood when you left me alone?} \\
\texttt{Claudia: That's not fair, Ty. I had to...} \\
\texttt{Tyler: What? Chase your dreams? Leave your past behind?} \\
\texttt{Claudia: I never left you behind, Tyler. You were always...} \\
\texttt{Tyler: Always what? An afterthought? A burden?} \\
\texttt{Claudia: No, never a burden. You were my responsibility.} \\
\texttt{Tyler: Responsibility? Is that why you're here now? After all these years?} \\
\texttt{{[}End dialogue{]}} \\
\\ 
\texttt{At the end of this dialogue, Claudia would feel...} \\
\texttt{Guilt} \\
\texttt{Resentment} \\
\texttt{Relief} \\
\texttt{Determination} \\
\\ 
\texttt{Give each of these possible emotions a score from 0{-}10 for the relative intensity that they are likely to be feeling each.} \\
\\ 
\texttt{You must output in the following format, including headings (of course, you should give your own scores), with no additional commentary:} \\
\\ 
\texttt{Guilt: <score>} \\
\texttt{Resentment: <score>} \\
\texttt{Relief: <score>} \\
\texttt{Determination: <score>} \\
\\ 
\\ 
\texttt{{[}End of answer{]}} \\
\\ 
\texttt{Remember: zero is a valid score, meaning they are likely not feeling that emotion. You must score at least one emotion > 0.} \\
\\ 
\texttt{Your answer:} \\
            }
        }
    }
     \\ \\

    \theutterance \stepcounter{utterance}  
    & \multicolumn{4}{p{0.6\linewidth}}{
        \cellcolor[rgb]{0.561,0.737,0.561}{
            \makecell[{{p{\linewidth}}}]{
                \texttt{\tiny{[P1$\rangle$GM]}}
                \texttt{Guilt: 2} \\
\texttt{Resentment: 9} \\
\texttt{Relief: 0} \\
\texttt{Determination: 1} \\
            }
        }
    }
    & & \\ \\

    \theutterance \stepcounter{utterance}  
    & & & \multicolumn{2}{p{0.3\linewidth}}{
        \cellcolor[rgb]{0.9,0.9,0.9}{
            \makecell[{{p{\linewidth}}}]{
                \texttt{\tiny{[GM$|$GM]}}
                \texttt{Guilt: 2} \\
\texttt{Resentment: 9} \\
\texttt{Relief: 0} \\
\texttt{Determination: 1} \\
            }
        }
    }
    & & \\ \\

    \theutterance \stepcounter{utterance}  
    & & & \multicolumn{2}{p{0.3\linewidth}}{
        \cellcolor[rgb]{0.9,0.9,0.9}{
            \makecell[{{p{\linewidth}}}]{
                \texttt{\tiny{[GM$|$GM]}}
                \texttt{Guilt: 8} \\
\texttt{Resentment: 4} \\
\texttt{Relief: 0} \\
\texttt{Determination: 6} \\
            }
        }
    }
    & & \\ \\

    \theutterance \stepcounter{utterance}  
    & & & \multicolumn{2}{p{0.3\linewidth}}{
        \cellcolor[rgb]{0.9,0.9,0.9}{
            \makecell[{{p{\linewidth}}}]{
                \texttt{\tiny{[GM$|$GM]}}
                \texttt{game\_result = LOSE} \\
            }
        }
    }
    & & \\ \\

    \theutterance \stepcounter{utterance}  
    & & \multicolumn{4}{p{0.6\linewidth}}{
        \cellcolor[rgb]{0.95,0.95,0.95}{
            \makecell[{{p{\linewidth}}}]{
                \texttt{\tiny{[P1$\leftrightarrow$P]}}
                \texttt{Post{-}episode Reflexion (not sent to GM):} \\
\texttt{New plan: My previous attempt failed because I misidentified the primary emotional driver. I scored "Resentment" at 9 and "Guilt" at 2, but the dialogue reveals Claudia is primarily defensive and remorseful, not resentful. Tyler is the one expressing resentment ("You always do," "Is that why you're here now?"). Claudia's lines ("That's not fair," "I never left you behind," "You were my responsibility") are attempts to justify her actions and alleviate her own guilt, not to attack Tyler. Therefore, the highest score should be **Guilt** (reflecting her struggle with the "responsibility" she abandoned), followed by **Determination** (her resolve to explain herself), while **Resentment** should be very low or zero as she is not attacking Tyler, and **Relief** should be zero due to the tension. Next time, I will assign the highest score to Guilt (e.g., 8{-}9), a moderate score to Determination (e.g., 3{-}4), and near{-}zero scores to Resentment and Relief.} \\
            }
        }
    }
    & \\ \\

\end{supertabular}
}

\end{document}
